\ifx\danslelivre\undefined
\documentclass[10pt]{../fichiers-configuration-latex/smfart}
\usepackage{ifthen}
\newboolean{english}
\newboolean{francais}
\setboolean{english}{false}
\setboolean{francais}{true}
\input{../fichiers-configuration-latex/paquets}
\selectlanguage{francais}
\input{../fichiers-configuration-latex/lemme-proposition-theoreme}
\input{../fichiers-configuration-latex/numerotation}
\input{../fichiers-configuration-latex/macros}

\externaldocument{../20-rigidite/rigidite}
\externaldocument{../04-Cohen-Gabber/Cohen-Gabber}

\title{XXI. Le théorème de finitude pour les coefficients non abéliens}
\author{Frédéric Déglise}

\begin{document}

\begin{center}
XXI. Le théorème de finitude pour les coefficients non abéliens\\
Frédéric Déglise
\end{center}

\else
\setboolean{english}{false}
\setboolean{francais}{true}
\selectlanguage{francais}
\chapter{Le théorème de finitude pour les coefficients
non abéliens}
\addtocontents{toc}{par Frédéric Déglise}
\begin{center}
Frédéric Déglise
\end{center}
\fi


\begin{flushright}
À la mémoire de mon oncle Olivier.
\end{flushright}

\section{Introduction}

Le but de l'exposé est de démontrer les théorèmes suivants,
qui généralisent le théorème d'Artin (cf \sga 4 {XIV}{th. 1.1})
dans le cas ensembliste et non abélien~:
\begin{theoreme}[Gabber] \label{21-ensembles}
Soit $f:Y \rightarrow X$ un morphisme de type fini
 entre schémas noethériens.

Pour tout faisceau constructible $F$ sur $Y_\et$, 
 le faisceau $f_*F$ est constructible.
\end{theoreme}

\begin{theoreme}[Gabber] \label{21-groupes}
Soit $f:Y \rightarrow X$ un morphisme de type fini
 entre schémas quasi-excellents.
 
Soit $\mathbb L$ un ensemble de nombres premiers inversibles sur $X$.

Pour tout faisceau constructible de groupes $F$ sur $Y_\et$ 
 de $\mathbb L$-torsion,
le faisceau $\R^1 f_*(F)$ sur $X_\et$ est constructible.
\end{theoreme}

\begin{theoreme}[Gabber] \label{21-groupes2}
Soit $X$ un schéma normal excellent,
 $Z \subset X$ une partie fermée de codimension supérieure à $2$.
Notons $j:U \rightarrow X$ l'immersion ouverte du complémentaire.

Pour tout groupe fini $G$,
 le faisceau $\R^1 j_*(G_U)$ est constructible.
\end{theoreme}
 
\begin{theoreme}[Gabber] \label{21-local}
Soit $A$ un anneau strictement local de dimension $2$.
On suppose que $A$ est normal, excellent, et on note $X'=\Spec(A)-\{\mathfrak{m}_A\}$
 son spectre épointé.

Alors, pour tout groupe fini $G$,
 l'ensemble $H^1(X';G)$ est fini.
\end{theoreme}

Le théorème \ref{21-ensembles} est prouvé dans la section \ref{21-sec-ensembles}.
Ce théorème est utilisé par les suivants dans le cas où $X$ est quasi-excellent.
Ce cas est beaucoup plus simple, comme nous le dégageons dans la démonstration.

Le théorème \ref{21-groupes} est réduit -- en trois étapes --
 au théorème \ref{21-groupes2} dans la section \ref{21-sec-groupes}.
Toutefois, le lecteur attentif notera que ce dernier théorème 
 n'est pas un simple cas particulier car il n'est pas nécessaire
 de faire d'hypothèse sur le cardinal du groupe $G$.

Le théorème \ref{21-groupes2} est réduit au théorème \ref{21-local}
 dans la section \ref{21-sec-groupes2}.
Cette réduction apparaît en \ref{21-sec-red-21-groupes2}
 et utilise deux lemmes qui ont été établis auparavant 
 (lemmes \ref{21-principeG} et \ref{21-bertini}).

Le dernier théorème est bien un cas particulier de \ref{21-groupes2}.
Toutefois, nous avons choisi de le dégager dans cette introduction
 à la fois comme une résultat important
  et comme un point clé.
Il est démontré dans la section \ref{21-sec-local}
 suivant un raisonnement par l'absurde qui utilise
 la méthode des ultrafiltres (voir \ref{21-sec-ultra-filtres} pour des rappels).

\paragraph*{Notations et conventions}
\begin{itemize}
\item Quand une topologie
 sur un schéma est sous-entendue, il s'agit de la topologie étale.
\item
Étant donné un ensemble $D$ (resp. un groupe $G$), on notera parfois $D$ (resp. $G$)
 pour le faisceau étale constant induit sur un schéma $X$ lorsque $X$ est clair
 d'après le contexte. Si l'on veut préciser $X$,
 on note ce faisceau $D_X$ (resp. $G_X$), suivant l'usage.
\item
Quand on parle de la normalisation d'un schéma $X$,
 il s'agit du morphisme canonique
$$
X'=\sqcup_{i \in I} X'_i \rightarrow X
$$
où $I$ désigne l'ensemble des composante irréductibles de $X$
 et $X'_i$ désigne le schéma normalisé de la composante irréductible de $X$ correspondant
 à $i$, munie de sa structure de sous-schéma réduit.
On dit aussi que $X'$ est le schéma normalisé associé à $X$.
\end{itemize}

\section{Image directe de faisceaux d'ensembles constructibles}
\label{21-sec-ensembles}

Dans le cas où $X$ est quasi-excellent,
la preuve est une application de résultats déjà connus
 (cf \sgacourt 4 {IX}). 
Nous commençons par exposer la démonstration dans ce
cas, puis dans le cas général. 
Toutefois, l'étape de réduction exposée dans la section qui suit
est valable dans les deux cas.

\subsection{Réduction du théorème}

On commence par réduire le théorème \ref{21-ensembles} à l'assertion suivante:
\begin{enumerate}
\item[\ppE]
Soit $D$ un ensemble fini et
$j:U \rightarrow X$ une immersion ouverte entre schémas noethériens.
Alors, le faisceau d'ensembles $j_*(D_U)$ est constructible.
\end{enumerate}
Considérons les hypothèses du théorème \ref{21-ensembles}.
D'après \sga 4 {IX} {2.14}, on peut trouver un monomorphisme
$$
F \rightarrow \prod_{i=1}^n \pi_{i*}(C_i)=Q
$$
pour des morphismes finis $\pi_i:Y_i \rightarrow Y$ et 
 des faisceaux constants finis $C_i$ sur $Y_i$.
Comme un sous-faisceau d'un faisceau constructible est constructible
(\sga 4 {IX} {2.9(ii)}),
il suffit de montrer que $f_*(Q)$ est constructible.
On est donc ramené au cas de $(f\pi_i)_*(C_i)$ pour tout $i$,
ce qui montre qu'on peut supposer $F=D_Y$ pour
un ensemble fini $D$.

Notons que dans ce cas, le théorème est local en $Y$.
En effet, si l'on se donne un recouvrement étale
$ \pi:W \rightarrow Y$ ($Y$ est noethérien),
 le morphisme d'adjonction 
$$
D_Y \rightarrow \pi_*\pi^*(D_Y)=\pi_*(D_W)
$$ 
est un monomorphisme. En lui appliquant $f_*$,
 on en déduit un monomorphisme
$$
f_*(D_Y) \rightarrow (f \pi)_*(D_W).
$$
Il suffit donc de montrer que le membre de droite est 
constructible (\sga 4 {IX} {prop. 2.9(ii)} à nouveau).
Notamment, on peut donc supposer que $Y$ est affine.

Alors, $f$ est séparé de type fini. 
On peut donc considérer une factorisation
$$
Y \xrightarrow j \bar X \xrightarrow{\bar f} X
$$
de $f$ telle que $j$ est une immersion ouverte et $\bar f$
un morphisme propre. Le résultat est connu pour $\bar f$ 
(cf \sga 4 {XIV} {th. 1.1}) donc on est réduit au cas 
de l'immersion ouverte $j$, c'est-à-dire à l'assertion
\ppE.

\remarque \label{21-rem-notes}
\begin{enumerate}
\item Si l'on suppose que $X$ est quasi-excellent,
 le schéma $\bar X$ qui apparaît dans la réduction ci-dessus
 est encore excellent puisque $\bar f$ est de type fini.
\item Dans cette réduction,
 on a vu que \ppE est locale en $U$.
\end{enumerate}

\subsection{Cas où $X$ est quasi-excellent}

Notons le lemme facile suivant:
\begin{lemme}\label{21-lm-normalisation}
Considérons un carré cartésien de schémas noethériens
$$
\xymatrix@=10pt{
U'\ar^{j'}[r]\ar_{q}[d] & X'\ar^{p}[d] \\
U\ar^j[r] & X
}
$$
tel que $j$ est une immersion ouverte
 et $p$ un morphisme fini surjectif.
Alors, pour tout ensemble fini $D$,
 si $l_*(D_V)$ est constructible, $j'_*(D_{U'})$ est constructible.
\end{lemme}
\begin{proof}
Par hypothèse, $q$ est surjectif. On en déduit que le morphisme
 d'adjonction
$$
D_U \rightarrow q_*q^*(D_U)=q_*(D_{V'})
$$
est un monomorphisme.
Appliquant $j_*$, on en déduit un monomorphisme 
$$j_*(D_U) \rightarrow p_*(j'_*(D_{U'})).
$$
Puisque $p$ est fini,
 $p_*$ préserve la constructibilité d'après \sga 4 {IX}{prop. 2.14(i)}.
Le lemme en résulte
puisqu'un sous-faisceau d'un faisceau d'ensembles constructible
 est constructible
 (\sga 4 {IX} {2.9(ii)}).
\end{proof}

Avant de passer à la preuve dans le cas général,
 notons que la démonstration du théorème \ref{21-ensembles}
 dans le cas où $X$ est quasi-excellent est plus simple.
Grâce à la remarque précédente, on se réduit à l'assertion
\ppE dans le cas où $X$ est quasi-excellent.
Puisque cette assertion est locale en $X$,
on peut supposer que $X$ est excellent.
Dès lors, la normalisation $p:X' \rightarrow X$ de $X$ est finie. 
Ainsi, le lemme précédent appliqué au carré cartésien évident
nous ramène au cas où $X$ est normal.

Seul le cas où $X$ est connexe nous intéresse. Alors, d'après
 \sga 4 {IX} {lem. 2.14.1}, $j_*(D_U)=D_X$, ce qui conclut.

\subsection{Cas général}

Considérons les hypothèse de l'assertion \ppE.

Notons que cette assertion est trivial pour $X=\emptyset$.
On peut donc raisonner par induction noethérienne sur $X$.
On suppose plus précisément que l'hypothèse d'induction suivante 
est vérifiée:
\begin{enumerate}
\item[($\mathscr H$)] Pour tout fermé strict $Z$ de $X$,
 pour tout morphisme fini $Z' \rightarrow Z$
 et pour toute immersion ouverte $l:V' \rightarrow Z'$,
 $l_*(D_{V'})$ est constructible.
\end{enumerate}
L'assertion étant locale en $X$,
 on peut supposer que $X$ est le spectre d'un anneau
 noethérien réduit $A$.
Utilisant le lemme \ref{21-lm-normalisation} -- en prenant 
pour $X'$ la somme disjointe des composantes irréductibles de $X$ --
on peut supposer aussi que $A$ est intègre.
On a déjà vu que \ppE est aussi locale en $U$
 (point ii de la remarque \ref{21-rem-notes}). 
On peut donc se ramener au cas où $U=\Spec(A_f)$ pour un élément
$f \in A$.

Soit $A'$ la clôture intégrale de $A$ dans son corps des fractions.
On pose $X'=\Spec(A')$, $Z=\Spec(A/(f))$
 et on considère le diagramme formé de carrés cartésiens:
$$
\xymatrix@R=15pt@C=25pt{
&& Z''\ar_h[d]\ar@/^14pt/^\pi[dd] \\
U'\ar^{j'}[r]\ar_q[d] & X'\ar_p[d] & Z'\ar_{i'}[l]\ar_r[d] \\
U\ar^j[r] & X & Z\ar_-i[l]
}
$$
tel que $i$ et $j$ sont les immersions évidentes,
 $p$ (resp. $h$) est la normalisation de $X$
 (resp. $Z'$).

\noindent {\it Première étape (fibres génériques de $p$)}: \\
Considérons les points génériques $z'_1,...,z'_n$ de $Z'$,
 et posons $z_r=p(z'_r)$. Bien que $p$ ne soit pas nécessairement
 fini, 
 $p^{-1}(z_r)$ est fini et l'extension résiduelle $\kappa(z'_r)/\kappa(z_r)$
 est finie (voir \cite[chap. V, th. 33.10]{Nag62}).
D'après le lemme \ref{21-lm-normalisation},
 on peut toujours remplacer $A$ par une extension finie
  $A \subset B \subset A'$.
L'hypothèse ($\mathscr H$) est en effet encore vérifiée pour $Y=\Spec(B)$.
Dès lors, on peut supposer que les conditions suivantes sont vérifiées:
\begin{enumerate}
\item[(h1)] Pour tout indice $r$, $p^{-1}(\{z_r\})=\{z'_r\}$.
\item[(h2)] Pour tout indice $r$, $\kappa(z'_r)/\kappa(z_r)$ est triviale.
\end{enumerate}
Notons $A_i$ (resp. $A'_i$) l'anneau localisé de $X$ en $z_i$ (resp. $X'$ en $z'_i$).
Alors, $A'_i$ est un anneau de valuation discrète.
Du fait que l'extension induite $A'_i/A_i$ est entière, on déduit que $A_i$
est de dimension $1$, ce qui implique que $z_i$ est un point générique du diviseur
$Z$ de $X$. Comme $q$ est surjectif, on déduit de (h1) que $z_1,...,z_n$ est l'ensemble
des points génériques de $Z$.

\noindent {\it Deuxième étape (restriction à un ouvert de $Z$)}: \\
Puisque trivialement $j^*j_*(D_U)=D_U$,
 il suffit de montrer que $i^*j_*(D_U)$ est constructible.
Notons les faits suivants:
\begin{enumerate}
\item $q$ surjectif: $D_U \rightarrow q_*q^*(D_U)=q_*(D_{U'})$
 est un monomorphisme.
\item $X'$ normal, $j'$ dominante: $j'_*(D_{U'})=D_{X'}$.\footnote{Comme
 on l'a déjà vu, c'est \sga 4 {IX} {lem. 2.14.1}}
\item $h$ surjectif: $D_{Z'} \rightarrow h_*h^*(D_{Z'})=h_*(D_{Z''})$
 est un monomorphisme.
\end{enumerate}
On déduit de {\bf i} et {\bf ii} un monomorphisme
$$
j_*(D_U) \rightarrow j_*q_*(D_{U'})=p_*(D_{X'}).
$$
Notons que $p$ est pro-fini. Le théorème de changement de base propre
 s'étend à ce cas, ce qui donne la relation: $i^*p_*=r_*{i'}^*$.
Si on applique $i^*$ au monomorphisme précédent,
 on déduit de cette relation et de {\bf iii} un monomorphisme composé:
$$
\sigma:i^*j_*(D_U) \rightarrow i^*p_*(D_{X'})=r_*(D_{Z'}) \rightarrow \pi_*(D_{Z''}).
$$
Considérons maintenant une immersion ouverte dense $l:V \rightarrow Z$
 ainsi que le carré cartésien suivant:
$$
\xymatrix@=15pt{
Z''\ar_\pi[d] & V''\ar_/-2pt/{l'}[l]\ar^{\pi_V}[d] \\
Z & V.\ar_/-2pt/{l}[l]
}
$$
On en déduit un diagramme commutatif de faisceaux d'ensembles:
$$
\xymatrix@R=15pt@C=20pt{
i^*j_*(D_U)\ar_\sigma[d]\ar^-\alpha[r] & l_*l^*i^*j_*(D_U)\ar^{l_*l^*\sigma}[d] \\
\pi_*(D_{Z''})\ar^-\beta[r] & l_*l^*\pi_*(D_{Z''}).
}
$$
Or le morphisme $\beta$ -- induit par le morphisme de coünité l'adjonction $(l^*,l_*)$ --
 est un isomorphisme.
En effet, $l$ étant une immersion ouverte $l^*\pi_*=\pi_{V*}{l'}^*$ par changement de base.
 On obtient donc l'identification: $l_*l^*\pi_*(D_{Z''})=\pi_*l'_*(D_{V''})$.
De plus, à travers cette identification, le morphisme $\beta$ est
l'image\footnote{On
 le vérifie facilement en revenant à la définition du morphisme de changement de base
 à l'aide des adjonctions $(l\*,l_*)$ et $({l'}^*,l'_*)$ et en utilisant que la composée suivante de morphismes unités/coünités est l'identité:
$$
l_* \rightarrow l_*l^*l_* \rightarrow l_*.
$$
}
par $\pi_*$ du morphisme de coünité pour l'adjonction $({l'}^*,l'_*)$:
$$
D_{Z''} \rightarrow l'_*l'(D_{Z''})=l'_*(D_{V''}).
$$
Ce dernier est un isomorphisme puisque $Z''$ est normal et $l'$ dense
 (voir \sga 4 {IX}{lem. 2.14.1}).
Puisque $\sigma$ est un monomorphisme,
 on en déduit que $\alpha$ est un monomorphisme. \\
Or, d'après ($\mathscr H$), le faisceau d'ensembles $l_*(D_V)$ est constructible.
Pour conclure, il suffit donc (d'après \sga 4 {IX} {2.9(ii)})
 de trouver un ouvert $V$ de $Z$ tel que
\begin{equation} \label{21-eq-immergées}
j_*(D_U)|_V=D_V.
\end{equation}

\noindent {\it Troisième étape (composantes immergées de $Z$)}: \\
Considérons la réunion $T$ des composantes immergées de $Z$ dans $X$.
Alors, l'ouvert $V=Z-T$ satisfait la relation \eqref{21-eq-immergées}.

Il s'agit de démontrer que pour tout point géométrique $\bar x$ de $V$,
 la fibre du morphisme canonique $D_V \rightarrow j_*(D_U)|_V$
 au point $\bar x$ est un isomorphisme.
Autrement dit, le morphisme induit
\begin{equation} \label{21-eq-immergéesP}
\pi_0(X_{(\bar x)}) \rightarrow \pi_0(X_{(\bar x)}-\{\bar x\})
\end{equation}
est un isomorphisme.

Soit $x$ le point de $X=\Spec A$ correspondant à $\bar x$.
Si $x$ n'est pas un point générique, 
puisque $V$ n'a pas de composante immergée,
$$
\mathrm{Prof}_x(A/(f)) \geq 1 \Rightarrow \mathrm{Prof}_x(A) \geq 2.
$$
D'après le théorème de Hartshorne (\sga 2 {III} {3.6}),
 le morphisme \eqref{21-eq-immergéesP} est donc un isomorphisme. \\
Supposons que $x$ est un point générique. D'après (h1), il existe
un indice $i$ tel que $x=z_i$. Or, d'après (h2), le morphisme
entier birationnel
$$
X'_{(z'_i)}=\Spec(A'_i) \rightarrow \Spec(A_i)=X_{(z_i)}
$$
est radiciel. C'est donc un homéomorphisme universel.
On en déduit que le morphisme $X'_{(\bar z'_i)} \rightarrow X_{(\bar z_i)}$
est encore un homéomorphisme, où $\bar z'_i$ est le point géométrique
correspondant à la clôture séparable de $\kappa(z_i)$ définie par $\bar z_i$.
Dès lors, \eqref{21-eq-immergéesP} est un isomorphisme 
puisque la propriété correspondante est vraie
 pour le schéma normal $X'_{(\bar z'_i)}$. Ceci conclut.

\section{Image directe dérivée de faisceaux de groupes constructibles}
\label{21-sec-groupes}

\subsection{Réduction au cas d'un faisceau constant.}

\label{21-constant}

\begin{lemme} \label{21-lm-sous-fx}
Soit $f:Y \rightarrow X$ une morphisme de type fini
 et $u:F \rightarrow F'$ un monomorphisme de faisceaux en groupes 
  constructibles sur $Y_\et$.
  
Alors, $\R^1 f_*(F')$ constructible implique
 $\R^1 f_*(F)$ constructible.
\end{lemme}
\begin{proof}
Posons $C=F'/F$ vu comme faisceau pointé, 
 constructible sur $Y_{\et}$ par hypothèse.
On considère la suite exacte de faisceaux pointés
 (cf \sga 4 {XII}{3.1})
$$
f_*(F') \rightarrow f_*(C) \rightarrow \R^1 f_*(F)
 \xrightarrow v \R^1 f_*(F').
$$
Supposant que $\R^1 f_*(F')$ est constructible,
on peut trouver une famille génératrice de sections locales 
$(e_1,...,e_n)$ de $\R^1 f_*(F')$.
Soit $\Phi_i$ le faisceau fibre de $v$ en $e_i$, 
défini par le diagramme cartésien de faisceaux sur $X_\et$~:
$$
\xymatrix@=10pt{
\Phi_i\ar[r]\ar[d] & V_i\ar^{e_i}[d] \\
\R^1 f_*(F)\ar^v[r] & \R^1 f_*(F').
}
$$
Utilisant le critère de constructibilité par les fibres et les
spécialisations (cf \sga 4 {IX}{prop. 2.13(ii)}),
on voit aisément qu'il suffit de montrer que le faisceau
$\Phi_i$ est constructible pour tout $i$.

Pour montrer cela, il suffit de se rappeler de l'interprétation de 
$\Phi_i$ en termes de $F$-objets tordus.
Soit $x$ un point de $X_\et$. 
Il existe un voisinage étale $V$ de $x$
dans $X$ et une section $e$ de $\Phi_i$ sur $V/X$.
Quitte à restreindre le voisinage $V$, $e$ provient d'un $F$-torseur
sur $V \times_Y X$. On peut alors tordre par $e$ la suite exacte 
ci-dessus restreinte à $V$, 
chaque faisceau étant muni d'une action de $f_*(F)$~:
$$
f_*(F')^e \rightarrow f_*(C)^e \rightarrow \R^1 f_*(F)^e
 \xrightarrow{v^e} \R^1 f_*(F')^e.
$$
La suite obtenue est encore exacte et 
$\Phi_i|_V$ s'identifie avec le noyau de $v^e$.
On en déduit $\Phi_i|_V \simeq f_*(C)^e/f_*(F')^e$.
Or ce faisceau est constructible d'après le théorème 
\ref{21-ensembles}.
\end{proof}

Considérons les hypothèses du théorème \ref{21-groupes}.
D'après \sga 4 {IX}{prop. 2.14}, on peut trouver un monomorphisme
de groupes
$$
F \rightarrow F'=\prod_{i=1}^r \pi_{i*}(G_i)
$$
pour des morphismes finis $\pi_i:U_i \rightarrow U$ 
 et des groupes finis $G_i$ pour $i=1,...,r$.
D'après le lemme précédent, on est réduit au cas de $F'$.
Puisque $\pi_{i*}$ est exact,
on est donc ramené au cas du morphisme $f \circ \pi_i$
et du faisceau constant sur $U_i$ de groupe $G_i$
 pour chaque indice $i$.

\subsection{Réduction au cas d'une immersion ouverte}

Le lemme clé dans cette étape de réduction est 
le suivant~:
\begin{lemme}
\label{21-devissage1}
Soit $G$ un groupe fini.
Considérons un diagramme commutatif
$$
\xymatrix@R=10pt@C=20pt{
 & {Y'}\ar^{g}[rd] & \\
Y \ar^{h}[ru]\ar|f[rr] && X
}
$$
de morphismes de type fini.
Les conditions suivantes sont vérifiées~:
\begin{enumerate}
\item[(i)] Si $h$ est surjectif, \\
$\R^1 f_*(G_Y)$ constructible implique $\R^1 g_*(G_{Y'})$ constructible.
\item[(ii)] Si $g$ est propre, \\
 $\R^1 h_*(G_Y)$ constructible implique $\R^1 f_*(G_Y)$ constructible.
\end{enumerate}
\end{lemme}
\begin{proof}
Rappelons que l'on dispose de la suite exacte de faisceaux
d'ensembles sur $X_\et$ (cf \sga 4 {XII}{prop. 3.2})~:
$$
* \rightarrow \R^1 g_*(h_* G_Y)
 \xrightarrow u \R^1 f_*(G_Y) \xrightarrow v g_*\R^1 h_*(G_Y).
$$

Considérons l'assertion (i). Puisque $h$ est surjectif,
 le morphisme d'adjonction 
$$
G_{Y'} \rightarrow h_* h^* G_{Y'}=h_* G_Y
$$ 
est un monomorphisme.
Appliquant le lemme \ref{21-lm-sous-fx}, il suffit de 
montrer que $\R^1 g_*(h_* G_Y)$ est constructible.
On peut alors conclure puisque le morphisme
$u:\R^1 g_*(h_* G_Y) \rightarrow \R^1 f_*(G_Y)$
est un monomorphisme.

Considérons maintenant l'assertion (ii).
Puisque $g$ est propre, 
 le théorème de changement de base propre \sga 4{XIV}{th. 1.1}
  conjugué avec le théorème \ref{21-ensembles} montre que
   la source de $u$ est constructible.
Par hypothèse et une nouvelle application du théorème \ref{21-ensembles},
le but de $v$ est constructible.

Il suffit alors de raisonner comme dans la démonstration de 
\ref{21-lm-sous-fx} sur les fibres du morphisme $v$ associées à une famille
finie de sections locales de $g_*\R^1 h_*(G_Y)$ qui est génératrice.
Chacune de ses fibres est localement vide ou isomorphe au faisceau tordu 
$\R^1 g_*\big((h_* G_Y)^e\big)$ pour une de ces sections locales $e$.
Comme ce faisceau est toujours constructible, on peut conclure.
\end{proof}

Considérons maintenant les hypothèses du théorème \ref{21-groupes},
dans le cas $F=G_Y$.
Puisque $Y$ est noethérien, 
 il existe un recouvrement Zariski $\pi:W \rightarrow Y$
  tel que $W$ est affine. 
D'après l'assertion (i) du lemme ci-dessus,
il suffit de montrer le théorème pour $f \circ \pi$.
On peut donc supposer que $Y$ est affine.

Le morphisme $f:Y \rightarrow X$ est alors quasi-projectif.
On peut donc considérer une factorisation 
$Y \xrightarrow j Y' \xrightarrow g X$ de $f$
tel que $g$ est projectif et $j$ est une immersion ouverte.
D'après l'assertion (ii) du lemme ci-dessus, nous sommes 
réduit au cas de l'immersion ouverte $j$.

\subsection{Réduction au théorème \ref{21-groupes2} (\emph{i.e.} la codimension $2$)}

Grâce aux deux étapes de réduction précédentes,
 nous sommes ramenés au cas d'une immersion ouverte
  $j:U \rightarrow X$ et d'un faisceau constant sur $U$
   de groupe $G$.
Dans cette étape de réduction, 
on considère la codimension du complémentaire $Z$ de $U$ dans $X$.

Pour montrer que $\R^1 j_*(G_U)$ est constructible
 on peut raisonner localement sur $X$.
On peut donc supposer que $X$ est excellent.
Considérons la normalisation $p:X' \rightarrow X$
de $X$ ainsi que le carré cartésien~:
$$
\xymatrix@=10pt{
U'\ar^{j'}[r]\ar_q[d] & X'\ar^p[d] \\
U\ar^j[r] & X.
}
$$
Puisque $q$ est surjectif, le morphisme d'adjonction
$G_U \rightarrow q_*q^*(G_U)$ est un monomorphisme.
D'après le lemme \ref{21-lm-sous-fx},
 il suffit donc de montrer que $\R^1 j_*(q_* G_{U'})$
 est constructible.
Or $p$ et $q$ étant finis,
 $\R^1 j_*(q_* G_{U'})=p_* \R^1 j'_*(G_{U'})$.
On peut donc supposer que $X$ est normal.

Dès lors, $X$ est somme disjointe de ses composantes
irréductibles, et on peut donc le supposer intègre.
Notons $K$ son corps des fonctions.
Pour une extension finie $L/K$,
 on peut considérer le schéma normalisé $X'$ de $X$ dans $L$,
 ainsi que le diagramme cartésien suivant~:
$$
\xymatrix@=10pt{
U'\ar^{j'}[r]\ar_q[d] & X'\ar^p[d] \\
U\ar^j[r] & X.
}
$$
Par adjonction, on obtient un morphisme canonique
$$
\phi^{(L)}_{U/X}:\R^1 j_*(G_U) \rightarrow p_* \R^1 j'_*(G_{U'})
$$
induit par le morphisme qui à un $G$-revêtement d'un
 schéma étale $V/U$ associe son pullback sur
  $V'=V \times_U U'$.
\begin{lemme}
\label{21-principeG}
Considérons les hypothèses et notations précédentes.
Alors, les conditions suivantes sont équivalentes~:
\begin{enumerate}
\item[(i)] $\R^1 j_*(G_U)$ est constructible.
\item[(ii)] Il existe une extension finie séparable $L/K$ 
telle que $\phi^{(L)}_{U/X}$ est trivial.
\end{enumerate}
\end{lemme}
\begin{proof}
$(ii) \Rightarrow (i)$~: Soit $L/K$ une extension finie
telle que $\phi^{(L)}_{U/X}$ est trivial.
Avec les notations qui précèdent,
on pose $C=q_*(G_{U'})/G_U$, faisceau sur $U_\et$ pointé 
 de manière évidente.
On peut alors former la suite exacte de faisceaux pointés
 (cf \sga 4 {XII}{3.2})~:
$$
j_*q_*(G_{U'}) \rightarrow j_*(C)
 \rightarrow \R^1 j_*(G_U)
  \xrightarrow{(1)} \R^1 j_*\left(q_*(G_{U'})\right).
$$
Notons que, puisque $p$ est fini,
$\R^1 j_*\left( q_*(G_{U'}) \right)=p_* \R^1 j'_*(G_{U'})$ et le morphisme
(1) s'identifie au morphisme $\phi^{(L)}_{U/X}$.
Par hypothèse, la suite exacte ci-dessus implique donc que
$\R^1 j_*(G)=j_*(C)/j_*q_*(G_{U'})$ ce qui conclut
 d'après le théorème \ref{21-ensembles}. \\
$(i) \Rightarrow (ii)$~: Considérons une famille $(e_1,...,e_n)$ 
de sections de $\R^1 j_*(G_U)$ qui soit génératrice.
Pour tout indice $i$, $e_i$ correspond à un revêtement
$P_i \rightarrow V_i \times_X U$ pour un schéma étale $V_i/X$.
On peut supposer que $P_i$ est connexe de corps des fonctions $L_i$.
Le morphisme $P_i \rightarrow X$ correspond donc à une extension
finie séparable $L_i/K$.
On considère la clôture normale $L$ d'une extension de $K$
 composée des $L_i/K$.
Par définition, $\phi^{(L)}_{U/X}(e_i)=*$ pour tout entier $i$
 et le résultat suit puisque $(e_1,...,e_n)$ est génératrice.
\end{proof}

\begin{remarque} \label{21-rem-principeG}
Considérons les notations qui précèdent le lemme.
En termes de $G$-revêtements, la trivialité du morphisme
$\phi^{(L)}_{U/X}$ s'interprète comme suit~:
\begin{enumerate}
\item[(i)] Pour tout point géométrique $\bar s$ de $Z$
  et pour tout $G$-revêtement 
   $\pi:P \rightarrow X_{(\bar s)}^h-Z_{(\bar s)}^h$,
le revêtement 
$\phi^{(L)}_{U/X,s}(P)$ est trivial.
\item[(ii)] Pour tout schéma étale $V/X$ et tout 
$G$-revêtement $\pi:P \rightarrow V-Z_V$,
 il existe un recouvrement étale $W/V'$ tel que $\pi|_{W-Z_W}$
 s'étend à $W$ (l'extension est alors unique puisque $X$ est normal).
\end{enumerate}
\end{remarque}

Dans le cas des immersions ouvertes, on peut renforcer
le lemme \ref{21-devissage1} comme suit~:
\begin{lemme}
\label{21-devissage2}
Considérons un diagramme commutatif
$$
\xymatrix@R=10pt@C=20pt{
 & V\ar^{k}[rd] & \\
U \ar^{h}[ru]\ar|j[rr] && X
}
$$
d'immersions ouvertes tel que $X$ est intègre normal et $h$ est dominante.
Alors, les conditions suivantes sont vérifiées~:
\begin{enumerate}
\item[(i)] Si $\R^1 j_*(G)$ est constructible, alors $\R^1 k_*(G)$ est constructible.
\item[(ii)] Si $\R^1 h_*(G)$ est constructible
 et pour tout morphisme fini surjectif $X' \rightarrow X$, $j'=j \times_X X'$,
  le faisceau $\R^1 j'_*(G)$ est constructible,
 alors $\R^1 j_*(G)$ est constructible.
\end{enumerate}
\end{lemme}
\begin{proof}
Remarquons que l'hypothèse sur $h$ et $V$ entraîne que 
 $h_*(G_U)=G_V$ (cf \sga 4 {IX}{lem. 2.14.1}).
 
L'assertion (i) résulte donc simplement du fait que
le morphisme canonique $\R^1 k_*(G_V) \rightarrow \R^1 j_*(G_U)$
est toujours un monomorphisme.

Considérons les hypothèses de l'assertion (ii).
D'après le lemme précédent appliqué à $h$,
 on peut trouver une extension finie $L/K$ telle que
  $\phi_{U/V}^{(L)}=*$.
Soit $X'$ le schéma normalisé de $X$ dans $L/K$,
 $k':V' \rightarrow X'$ le pullback de $k$ sur $X'$.
Appliquant le lemme précédent à $k'$, on peut trouver
une extension finie $E/L$ telle que $\phi_{V'/X'}^{(E)}=*$.

On note $X''$ le normalisé de $X$ dans $E/K$,
 $X'' \xrightarrow{p'} X' \xrightarrow p X$ les morphismes
  canoniques. On note $h'$, $j'$, $k'$ 
   (resp. $h''$, $j''$, $k''$) 
     les pullback respectifs
    de $h$, $j$, $k$ sur $X'$ (resp. $X''$).
On peut alors conclure grâce au diagramme commutatif
 suivant~:
$$
\xymatrix{
& \R^1j_*(G)\ar[r]\ar^{\phi^{(L)}_{U/X}}[d]\ar@{-->}[ld]
 & k_*\R^1 h_*(G)\ar^{k_*\phi^{(L)}_{U/V}}[d] \\
p_*\R^1k'_*(G)\ar[r]\ar_{p_*\phi^{(E)}_{V'/X'}}[d]
 & p_* \R^1j'_*(G)\ar[r]\ar^{p_*\phi^{(E)}_{U'/X'}}[d]
 & (pk')_*\R^1 h'_*(G) \\
(pp')_*\R^1k''_*(G)\ar[r] & (pp')_* \R^1j''_*(G).
 &
}
$$
La flèche pointillée existe du fait de l'exactitude de
la suite horizontale du milieu et de $\phi_{U/V}^{(L)}=*$.
On conclut puisque $\phi^{(E)}_{U/X}=
 (p_* \phi^{(E)}_{U'/X'}) \circ \phi_{U/X}^{(L)}$.
\end{proof}

\begin{lemme}
Soit $X$ un schéma régulier
 et $Z$ un sous-schéma fermé régulier de $X$.
Notons $j:U \rightarrow X$ l'immersion ouverte complémentaire.

Soit $n$ l'ordre de $G$.
Alors, si $n$ est inversible sur $X$,
 $\R^1 j_*(G_U)$ est constructible.
\end{lemme}
\begin{proof}
L'assertion est évidemment locale en $X_\et$.
On peut donc supposer que $X$ est local strictement hensélien
 et régulier, spectre d'un anneau noté $A$.
 
Si $Z$ est de codimension supérieure à $2$ dans $X$,
 il résulte du théorème de pureté de Zariski-Nagata
  (cf \sga 1 {X} {3.1, 3.3}) que $\R^1 j_*(G)=*$.

En codimension $1$, $Z$ admet un paramètre régulier $f \in A$.
On considère le schéma strictement local 
$$
X'=\Spec(A[t]/t^n-f).
$$
Le lemme d'Abhyankar absolu (cf \sga 1 {XIII}{prop. 5.2})
montre alors précisément que
 pour tout revêtement $E \xrightarrow \pi U$
  principal galoisien de groupe $G$,
le revêtement $\pi \times_X X':E' \rightarrow U'$ se prolonge
 à $X'$. 
Il est donc trivial et le lemme \ref{21-principeG} 
 accompagné de la remarque \ref{21-rem-principeG}
  permet de conclure.
\end{proof}

Revenons au cas général d'une immersion ouverte
 $j:U \rightarrow X$ de fermé complémentaire $Z$,
  $X$ étant supposé normal. Supposons que l'ordre de $G$
   est inversible sur $X$.

Soit $T$ la réunion des lieux singuliers de $X$ et $Z$.
Posons $V=X-T$ et $W=X-(Z \cup T)$ et
 considérons les immersions ouvertes correspondantes~:
$$
\xymatrix@R=-2pt@C=20pt{
& V\ar^k[rd] & \\
W\ar^h[ru]\ar_l[rd]\ar|\nu[rr] && X \\
& U\ar_j[ru] & 
}
$$
D'après le lemme précédent, $\R^1 h_*(G)$ est constructible.
D'après le lemme \ref{21-devissage2},
 on est donc ramené à prouver que pour tout morphisme fini 
  surjectif $X' \rightarrow X$, $k'=k \times_X X'$,
   le faisceau $\R^1 k'_*(G)$ est constructible.
Cette dernière assertion est bien impliqué par le théorème 
\ref{21-groupes2}.

\section{Cas de codimension $2$ sans hypothèse sur la torsion}
\label{21-sec-groupes2}

\subsection{Résolution des singularités}

\begin{lemme} \label{21-trivial-kernel}
Soit $X$ un schéma normal connexe,
$Z \subset X$ une partie fermée de codimension $2$
 et $j:U \rightarrow X$ l'immersion ouverte complémentaire.

Supposons $\mathrm{dim}(Z)>0$.
Alors, il existe une partie fermée $T \subset Z$ 
de codimension $3$ dans $X$ telle que pour tous points géométriques 
$\bar s$ et $\bar t$ de $Z-T$
 et toute spécialisation
  $\eta:X_{(\bar t)}^h \rightarrow X_{(\bar s)}^h$,
le noyau du morphisme de spécialisation
$$
\eta^*:\R^1 j_*(G)_{\bar s} \rightarrow \R^1 j_*(G)_{\bar t}
$$
est trivial.
\end{lemme}
\begin{proof}
Cette assertion ne dépend que d'un voisinage ouvert de $Z$ dans $X$.
Quitte à enlever une partie nulle part dense de $Z$,
on peut donc supposer que le lieu singulier de $X$ est inclus dans $Z$.
De même, on peut supposer que les composantes irréductibles de $Z$ sont 
disjointes, et de là, que $Z$ est irréductible.

D'après \cite{D2DS@Lipman},
 on peut résoudre la singularité de $X$ au point générique de $Z$
  par une suite d'éclatements et de normalisations.
Donc, quitte à retirer de nouveau une partie fermée nulle part dense de $Z$, 
on peut supposer qu'il existe un diagramme formé
de carrés cartésiens
$$
\xymatrix@=22pt{
Z'\ar[r]\ar[d]
 & X'\ar_p[d] & U'\ar^q[d]\ar|f[ld]\ar_/-2pt/{j'}[l] \\
Z\ar[r] & X & U\ar_/-5pt/j[l]
}
$$
tel que $X'$ est régulier, $Z$' est un diviseur dans $X'$,
toute composante irréductible de $Z'$ domine 
 $Z$ et $q$ est un isomorphisme.

On en déduit donc $\R^1 j_*(G)=\R^1 f_*(G)$.
Comme $X'$ est régulier et $U'$ dominant dans $X'$,
on obtient $j'_*(G_{U'})=G_{X'}$ (cf \sga 4 {IX}{lem. 2.14.1}).
On obtient donc un monomorphisme canonique 
$\rho:\R^1 p_*(G) \rightarrow \R^1 f_*(G)$.

Considérant les notations du lemme (où l'on a supposé $T=\emptyset$), 
on obtient donc un diagramme commutatif d'ensembles pointés
$$
\xymatrix{
\R^1 p_*(G)_{\bar s}\ar^{\rho_s}[r]\ar[d]
 & \R^1 f_*(G)_{\bar s}\ar^{\eta^*}[d] \\
\R^1 p_*(G)_{\bar t}\ar[r]
 & \R^1 f_*(G)_{\bar t}.
}
$$
D'après le théorème de changement de base propre appliqué à $p$, 
 la composée $\eta^* \circ \rho_s$ est un monomorphisme.
Il nous suffit donc de vérifier que le noyau de $\eta^*$
 est inclus dans l'image de $\rho_s$.

Quitte à tirer la situation sur $X_{(\bar s)}^h$,
on peut supposer que $X=X_{(\bar s)}^h$ pour simplifier les
notations.
On se donne donc un $G$-revêtement principal connexe
 $P \rightarrow U'$ qui est trivial
  sur $U' \times_X X_{(\bar t)}^h$.
On peut supposer que $P$ est connexe.
Si $F$ désigne le corps des fonctions de $X'$, 
 ce revêtement correspond à une extension séparable $E/F$.
Soit $\bar P$ la clôture normale de $X'$ dans $E/F$.
Bien sûr, $\bar P \times_{X'} U'=P$ puisque $P$ est 
  la clôture normale de $U'$ dans $E/F$.
Donc $\bar P/X'$ est non ramifié au-dessus de $U'$.
Du fait que $P \times_X X_{(\bar t)}^h$ est trivial,
 il suit que $\bar P$ est non ramifié en tous points 
  génériques  de $Z'$
   (puisque ceux-ci dominent $Z_{(\bar t)}^h$).
Autrement dit, $\bar P$ est non ramifié en codimension $1$.
Mais d'après le théorème de Zariski-Nagata (cf \sga 1 {X}{th. 3.1}),
 $\bar P$ est non ramifié en codimension supérieure à $2$
  car $X'$ est régulier.
Ainsi, $\bar P$ est le revêtement étale cherché.
\end{proof}

\subsection{Un argument « à la Lefschetz »}

Pour cet argument, nous utiliserons le lemme suivant,
 qui est une application des résultats liés à « la méthode de
 Lefschetz » de \sga 2 {X}{\S 2}.

\begin{lemme}\label{21-ref-deg-connexe}
Soit $X$ un schéma normal excellent connexe,
 $D$ un diviseur de Cartier effectif connexe dans $X$,
  et $Z \subset D$ une partie fermée de codimension supérieure à $2$.
Alors, $D-Z$ est connexe.
\end{lemme}

\begin{proof}
Il suffit de montrer que pour tout point $s$ de $X$, le schéma 
 $(D-Z) \times_X X_{(s)}$ est connexe.
On peut même supposer que $s$ est un point générique de $Z$.
On peut donc supposer que $X$ est local, de dimension supérieure à $3$
 et que $Z$ est son point fermé.
Considérons le complété $\hat X$ du schéma local $X$.
Puisque $X$ est excellent, $\hat X$ est encore normal.
Puisque le morphisme $\hat X \rightarrow X$ est un épimorphisme universel,
 il suffit de montrer que $(D-Z) \times_X \hat X$ est connexe.
On peut donc supposer en outre que $X$ est complet.
 
Alors, le spectre épointé $X'=X-Z$ est normal, connexe et 
de dimension supérieure à $2$. 
Il résulte du {\it critère de normalité de Serre}
(cf \cite[23.8]{CRT@Matsumura})
 que pour tout point fermé $x$ de $X'$,
$$\prof(\mathscr O_{X',x}) \geq 2.$$
Dès lors, d'après \sga 2 X {ex. 2.1} (voir aussi plus directement
 \sga 2 {IX}{prop. 1.4}), on obtient un isomorphisme canonique
$$
\Gamma(X',\mathscr O) \simeq
 \Gamma\left(\widehat{X'}^D,\mathscr O\right)
$$
où $\widehat{X'}^D$ désigne le complété formel de $X'$ le long de $D$,
 ce qui conclut.
\end{proof}

\begin{lemme} \label{21-bertini}
Soit $X$ un schéma normal excellent,
 $D$ un diviseur principal dans $X$
  et $Z \subset D$ une partie fermée non vide de codimension 
   supérieure à $2$.
On pose $U=X-Z$, $V=D-Z$.

Considérons le carré cartésien formé des immersions
évidentes
$$
\xymatrix@=10pt{
V\ar^{j'}[r]\ar_{i_U}[d] & D\ar^i[d] \\
U\ar^j[r] & X.
}
$$
Alors, le morphisme de changement de base associé
$$
i^* \R^1 j_*(G) \rightarrow \R^1 j'_*(G)
$$
est un monomorphisme.
\end{lemme}
\begin{proof}
On peut supposer que $X$ et $D$ sont locaux strictement
henséliens. On doit montrer que le morphisme 
de restriction
$$
H^1(U;G) \xrightarrow{i_U^*} H^1(V;G)
$$
est injectif.

Remarquons que le lemme précédent nous montre déjà que $V$
 est connexe.
Soit $P$ et $P'$ deux $G$-torseurs sur $U$ qui coïncident
sur $V$. On considère le faisceau 
$L=\underline{\mathrm{Isom}}_G(P,P')$
des $G$-isomorphismes de $P$ dans $P'$ sur $U_\et$.
on doit montrer qu'il admet une section sur $U$.

Comme $L$ est localement constant constructible,
 il est représentable par un $U$-schéma étale fini noté $U'$.
Posons $V'=U' \times_U V$. 
Par hypothèse, $V'/V$ admet une section.

On va montrer que pour toute composante connexe $U'_0$ de $U'$
telle que $U'_0 \times_U V/V$ admet une section, il
existe une section de $U'_0/U$ ce qui suffira pour conclure.

Quitte à remplacer $U'$ par $U'_0$, on peut supposer 
pour montrer cela que $U'$ est connexe.
Les corps de fonctions de $U'$ et $U$ 
définissent une extension finie séparable $L/K$.
Soit $X'$ le schéma normalisé de $X$ dans $L/K$. 
On pose encore $D'=X' \times_X D$ et $Z'=X' \times_X Z$.

Notons que $X'$ est normal excellent et connexe. 
Dès lors, le lemme précédent implique que
  $V'=D'-Z'$ est connexe.
Par hypothèse, le $V$-schéma étale $V'$ admet une section, 
 donc $V'=V$. 
Il en résulte que le revêtement étale $U'/U$ est de degré $1$ 
au-dessus de $V$, donc $U'=U$.
\end{proof}

\subsection{Réduction au théorème \ref{21-local}}
\label{21-sec-red-21-groupes2}

On peut supposer que $X$ est affine.
Il alors limite projective de schémas affines de type fini sur
$\Spec(\ZZ)$. Comme il suffit de démontrer le théorème \ref{21-groupes2}
pour chacun des termes de cette limite projective, on est réduit
au cas où $X$ est de dimension finie.

On raisonne par induction sur la dimension de $X$.
Le cas où $X$ est de dimension $2$ résulte du théorème \ref{21-local}.

On se place dans le cas où $\dim(X)>2$.

Si $\codim(Z)>2$,
 on peut trouver un diviseur principal connexe $D \xrightarrow i X$
  qui contient $Z$.
Soit $j':D-Z \rightarrow D$ le morphisme induit. \\
D'après le lemme \ref{21-bertini}, 
on obtient un monomorphisme $i^* \R^1 j_*(G) \rightarrow \R^1 j'_*(G)$.
Par hypothèse de récurrence, $\R^1 j'_*(G)$ est constructible.
On peut donc conclure car $\R^1 j_*(G)=i_*i^*\R^1 j_*(G)$.

Plaçons nous dans le cas critique où $\codim(Z)=2$,
 ce qui entraîne $\dim(Z)>0$.

D'après hypothèse d'induction,
le théorème est connu pour le schéma semi-localisé de $X$
aux points génériques de $Z$ qui sont de codimension $2$
dans $X$. Il existe donc, d'après le lemme \ref{21-principeG},
une extension finie $L$ du corps des fonctions $K$ de $X$ 
telle que 
pour tout point générique $\eta$ de $Z$ de codimension $2$ dans $X$,
$\phi_{U/X,\eta}^{(L)}$ est trivial.

Considérons $X'$ la normalisation de $X$ dans $L/K$ et
$p:X' \rightarrow X$ sa projection.
On pose $j'=j \times_X X'$ et $Z'=Z \times_X X'$.
D'après le lemme \ref{21-trivial-kernel},
il existe une partie fermée $T' \subset Z'$ de codimension supérieure à $3$
dans $X'$ tel que les noyaux des flèches de spécialisations de
$\R^1 j'_*(G)$ aux points de $Z'-T'$ soient triviaux.

Soit $T$ la réunion des composantes irréductibles de $Z$ de codimension 
supérieure à $3$ et du fermé $p(T')$ dans $Z$.
Considérons un point géométrique $\bar s$ de $Z-T$.
Il existe un point générique géométrique $\bar t$ de $Z-T$ et
une spécialisation $\eta:X_{(\bar t)}^h \rightarrow X_{(\bar s)}^h$.
Considérant les fibres de $\phi^{(L)}_{U/X}$, on obtient le diagramme
suivant~:
$$
\xymatrix@R=20pt@C=34pt{
\R^1 j_*(G)_{\bar s}\ar^-{\phi^{(L)}_{U/X,\bar s}}[r]\ar[d]
 & [p_* \R^1 j'_*(G)]_{\bar s}\ar^{\eta^*}[d] \\
\R^1 j_*(G)_{\bar t}\ar^-{\phi^{(L)}_{U/X,\bar t}}[r]
 & [p_* \R^1 j'_*(G)]_{\bar t}.
}
$$
D'après le choix de $L/K$,
la composée $\eta^* \circ \phi^{(L)}_{U/X,\bar s}$ est triviale.
Par ailleurs, puisque $p$ est fini, $\eta^*$ a un noyau trivial.
On en déduit que $\phi_{U/X,\bar s}^{(L)}$ est trivial.
D'après le lemme \ref{21-principeG}, $\R^1 h_*(G)$ est constructible
 pour l'immersion ouverte $h:X-Z \rightarrow X-T$.
D'après le lemme \ref{21-devissage2},
 on est donc réduit à montrer que pour tout morphisme fini surjectif 
  $X' \rightarrow X$, $\R^1 k'_*(G)$ est constructible
   pour l'immersion ouverte $k':X'-T' \rightarrow X'$.

Or on peut trouver un diviseur principal $D \xrightarrow i X'$
 qui contient $T'$.
On pose $k''=k' \times_{X'} D$.

D'après le lemme \ref{21-bertini}, 
on obtient un monomorphisme $$i^* \R^1 k'_*(G) \rightarrow \R^1 k''_*(G).$$
On peut donc à nouveau conclure d'après l'hypothèse d'induction appliquée
à $k''$ et du fait que $\R^1 k'_*(G)=i_*i^* \R^1 k'_*(G)$.

\section{Revêtements principaux d'une surface strictement locale épointée}
\label{21-sec-local}

\subsection{Mise en place}

D'après le théorème de rigidité de Gabber (cf \refext{XX}{20-prop-rig}),
 on peut supposer que $A$ est complet.
Soit $X=\Spec(A)$ et $X'=\Spec(A)-\{\mathfrak{m}_A\}$.

D'après le théorème de Cohen-Gabber (cf \refext{IV}{4-Cohen-Gabber égale}),
 il existe un sous-anneau régulier $R \subset A$
 tel que $A/R$ est finie génériquement étale.
On note $m$ le degré générique de $A/R$
Notons tout de suite le fait suivant qui résulte de l'algèbre
commutative standard~:
\begin{lemme} \label{21-facile}
Soit $R$ un anneau local régulier de dimension $2$,
 $A/R$ une algèbre finie tel que $A$ est normal.
Soit $m$ le degré générique de $A/R$.

Alors $A$ est un $R$-module libre de rang $m$.
\end{lemme}
En effet, puisque $R$ est local régulier de dimension $2$,
 par définition $\prof(R)=2$.
Par ailleurs, puisque $A$ est local normal de dimension $2$,
 il résulte du critère de Serre que $\prof(A)=2$ (cf \cite[ex.
17.3]{CRT@Matsumura}).
D'après le théorème d'Auslander-Buchsbaum (cf \cite[th.
19.1]{CRT@Matsumura}),
$$\prof(A)+\mathrm{dim\ proj}(A)=\prof(R).$$ Il en résulte que 
le $R$-module $A$
 est nécessairement projectif, donc libre.

Faisant abstraction du groupe $G$,
 on fixe un entier $n>0$ et on montre que l'ensemble
  des revêtements étales de $X'$ de degré $n$ est fini.

On raisonne par l'absurde.
Considérons une suite $(P'_i \rightarrow X')_{i \in \NN}$
 de revêtements étales de degré $n$ telle que pour
  tout $i \neq j$, $P'_i$ est non $X'$-isomorphe à $P'_j$.

Soit $K$ le corps des fractions de $A$.
Pour tout entier $i$, $P'_i/X'$ correspond
 à une extension finie séparable $L_i/K$.
On note $B_i$ la clôture intégrale de $A$ dans $L_i$,
 $P_i=\Spec(B_i)$.
Remarquons par ailleurs que d'après le lemme précédent,
 $B_i/R$ est libre de rang $nm$.

\subsubsection{Questions de discriminant}

Rappelons la définition suivante~:
\begin{definition}
Soit $B/A$ une algèbre finie libre de rang $n$.
Soit $\mathscr B=(e_i)_{1 \leq i \leq n}$ une base de $B/A$.
Le déterminant de la matrice 
 $\left(\mathrm{Tr}_{B/A}(e_ie_j)\right)_{1 \leq i,j \leq n}$
est appelé le discriminant de $B/A$ relativement à $\mathscr B$.
Sa classe dans le monoïde multiplicatif $A/(A^\times)^2$
 est indépendante de $\mathscr B$.
On la note $\disc_{B/A}$.
\end{definition}
Par abus, on considérera la classe $\disc_{B/A}$ 
 comme un élément de $A$.
Rappelons que $B/A$ est étale si et seulement si
 $\disc_{B/A}$ est inversible dans $A$.
Par la suite, nous aurons besoin de la formule suivante
(cf \cite[2.1.4]{lm@Ramero})~:
Soit $B/A$ et $C/B$ deux algèbres finies libres.
Soit $n$ le rang de $C/B$.
Alors,
$$
\disc_{C/A}=\disc_{B/A}^n.N_{B/A}(\disc_{C/B}).
$$

Revenant à la situation du numéro précédent,
 on considère un idéal $\p$ de hauteur $1$ de $R$.
Soit $A_\p$ (resp. $B_{i,\p}$) l'anneau semi-localisé de $A$
(resp. $B_i$) correspondant à la fibre au-dessus de $\p$.

Notons que $A_\p$ est normal de dimension $1$.
Il résulte du lemme \ref{21-facile} que $B_{i,\p}/A_\p$
 est libre de rang $n$.
D'après la formule rappelée précédemment,
$$
\disc_{B_{i,\p}/R_\p}
=\disc_{A_\p/R_\p}^n.N_{A_\p/R_\p}(\disc_{B_{i,\p}/A_\p}).
$$
Or $A/R$ (resp. $B_i/R$) est génériquement étale
 et $B_{i,\p}/A_\p$ est étale.
On déduit de la relation précédente que l'élément 
$(\disc_{B_{i}/R})(\disc_{A/R}^n)^{-1}$
de $K^\times$ appartient à $R^\times_\p$. 
Comme ceci est valable pour tout $\p$
et que $R$ est normal, on en déduit~:
\begin{equation} \label{21-disc}
\frac{\disc_{B_{i}/R}}{\disc_{A/R}^n} \in R^\times
\end{equation}
 
\subsection{Lemme clé}

\subsubsection{Ultraproduits} \label{21-sec-ultra-filtres}

\begin{definition}
Soit $I$ un ensemble.
Un ultrafiltre $\F$ sur $I$ est la donnée d'un ensemble
 de parties de $I$ ordonné par inclusion vérifiant les propriétés
 suivantes~:
\begin{enumerate}
\item $\forall F \in \F, \ \forall G \in P(I), \ F \subset G \Rightarrow \ G \in \F$.
\item $\forall F,G \in \F, \ F \cap G \in \F$.
\item $\forall F \in P(I)$, $F \in \F$ ou bien $I\backslash F \in \F$.
\item $X \in \F$ et $\emptyset \notin \F$.
\end{enumerate}
\end{definition}

\begin{exemple}
Soit $a$ un élément de $I$. Alors, l'ensemble des parties de $I$
 contenant $a$ est un ultrafiltre $\F$ de $I$.
Dans ce cas, on dit que $\F$ est \textit{principal}.
\end{exemple}

D'après le lemme de Zorn, il existe des ultrafiltres
 non principaux sur un ensemble infini $I$.

\begin{definition}
Soit $\F$ un ultrafiltre sur un ensemble $I$
 et $\C$ une catégorie admettant des limites inductives filtrantes
  et des produits.

Soit $(X_i)_{i \in I}$ une famille d'objets de $\C$.
Le système inductif $(\prod_{i \in F} X_i)_{F \in \F}$ est filtrant.
On définit l'ultraproduit de $(X_i)_{i \in I}$
 suivant $\F$ comme la limite inductive de ce système~:
$$
\prod_{i\in I/\F} X_i
 =\colim_{F \in \F} \left( \prod_{i \in I} X_i \right).
$$

Si $(X_i)_i$ est la famille constante de valeur un objet $X$,
 on note $X^{I/\F}$ son ultraproduit, appelé l'ultrapuissance
  de $X$ suivant $\F$.
On dispose toujours de l'application diagonale
 $X \rightarrow X^{I/\F}$.
\end{definition}

Nous utiliserons cette notion dans le cas des anneaux ou des modules.
Nous aurons besoin du lemme élémentaire suivant~:
\begin{lemme} \label{21-lm-ultra-prod}
Soit $I$ un ensemble et $\F$ un ultrafiltre sur $I$.

Considérons une famille $(A_i)_{i \in I}$ d'anneaux.
On pose $A_\infty=\prod_{i \in I/\F} A_i$.
\begin{enumerate}
\item[(i)] Si pour tout $i \in I$,
 $A_i$ est intègre (resp. un corps),
  il en est de même de $A_\infty$.
\item[(ii)]  Si pour tout $i \in I$,
 $A_i$ est local d'idéal maximal $\mathfrak{m}_i$,
  $A_\infty$ est local d'idéal maximal $\prod_{i \in I/\F} \mathfrak{m}_i$.
\end{enumerate}

Considérons une famille d'algèbres $(B_i/A_i)_{i \in I}$,
 $B_\infty/A_\infty$ son ultraproduit suivant $\F$~:
\begin{enumerate}
\item[(iii)] Si pour tout $i \in I$,
 $B_i/A_i$ est locale (resp. libre de rang $m$),
  il en est de même de $B_\infty/A_\infty$.
\end{enumerate} 

Considérons un anneau $A$, et $A_\infty=A^{I/F}$ son ultrapuissance.
\begin{enumerate}
\item[(iv)] Si $M$ est un $A$-module de présentation finie, \\
 $M \otimes_A A^{I/\F}=M^{I/\F}$.
\item[(v)] Si $A$ est cohérent, l'application diagonale
 $A \rightarrow A^{I/\F}$ est plate.
\end{enumerate}
\end{lemme}

\subsubsection{Énoncé et conséquence du lemme clé}

Revenons à la situation qui nous occupe.
Soit $\F$ un ultrafiltre non principal sur $\NN$.
On note $B_\infty$ (resp. $A_\infty$, $R_\infty$)
l'anneau local obtenu par ultraproduit suivant $\F$
de $(B_i)_{i \in \NN}$ (resp. $A$, $R$).
Soit $\widehat{B_\infty}$, $\widehat{A_\infty}$,
 $\widehat{R_\infty}$ leurs complétions respectives.
On en déduit les tours d'anneaux locaux suivantes~:
$$
\xymatrix@R=12pt@C=30pt{
B_i\ar[r] & B_\infty\ar[r]
  & \widehat{B_\infty}
 \\
A\ar[u]\ar[r] & A_\infty\ar^\varphi[u]\ar[r]
  & \widehat{A_\infty}\ar^{\hat \varphi}[u] \\
R\ar[u]\ar[r] & R_\infty\ar[u]\ar[r]
   & \widehat{R_\infty}.\ar[u]
}
$$
Considérant le spectre des anneaux locaux sur les deux premières lignes, 
on obtient le diagramme commutatif suivant~:
$$
\xymatrix@R=12pt@C=30pt{
P_i\ar[d] & P_\infty\ar[l]\ar^p[d]
 & \hat P_\infty\ar[l]\ar^{\hat p}[d] \\
X & X_\infty\ar[l] & \hat X_\infty\ar[l]
}
$$
Considérant encore le spectre épointé des anneaux locaux $A$, $A_\infty$
 et $\widehat{A_\infty}$, on obtient par changement de base le diagramme
 suivant~:
$$
\xymatrix@R=12pt@C=30pt{
P_i'\ar[d] & P_\infty'\ar[l]\ar^{p'}[d]
 & \hat P_\infty'\ar[l]\ar^{\hat p'}[d] \\
X' & X_\infty'\ar[l] & \hat X_\infty'\ar[l]
}
$$\begin{lemme}
\begin{enumerate}
\item[(1)] Le schéma $\hat X_\infty$ est noethérien.
\item[(2)] Le morphisme $\hat X_\infty \rightarrow X$ est plat
 à fibre géométriques régulières.
\item[(3)] Le morphisme $\hat X_\infty \rightarrow X_\infty$
 est une immersion fermée surjective sur les points fermés
  de $X'_\infty$.
\item[(4)] Le morphisme fini
  $\hat p:\hat P_\infty \rightarrow \hat X_\infty$
   est étale au-dessus de $\hat X'_\infty$.
\item[(5)] Le morphisme fini
  $p:P_\infty \rightarrow X_\infty$
   est étale au-dessus de $X'_\infty$.
\end{enumerate}
\end{lemme}
Montrons pourquoi ce lemme permet de terminer la démonstration.

D'après la variante C2 du théorème de changement de base lisse
 (cf \refext{XX}{20rigfaible}, cas (ii))
 appliquée au morphisme $X \rightarrow \hat X_\infty$
  et à l'ouvert $\hat X_\infty'$,
  il existe un revêtement étale $Q \rightarrow X'$ tel que
   $\hat P'_\infty \simeq Q \times_{X'} \hat X'_\infty$.

D'après le théorème de rigidité de Gabber (cf \cite{})
appliqué à $A_\infty$, le morphisme
$$
H^1(X'_\infty,G) \rightarrow H^1(\hat X'_\infty,G)
$$
est un isomorphisme. 
Donc $P'_\infty  \simeq Q \times_{X'} X'_\infty$.

Rappelons que le morphisme $p'$
est la limite projective cofiltrante
suivant les éléments $F$ de l'ultrafiltre $\F$
des morphismes canoniques
$$
p'_F:\bigsqcup_{i \in F} P'_i \rightarrow \bigsqcup_{i \in F} X'.
$$

Mais alors, les données étant de présentation finie sur $X'$,
 il existe un élément $F$ de $\F$ tel que
  $\bigsqcup_{i \in F} P'_i \simeq \bigsqcup_{i \in F} Q$.
Ceci constitue une contradiction puisque, $\F$ étant non principal,
 $F$ contient au moins deux éléments.

\subsubsection{Preuve du lemme clé}~

\medskip

\noindent {\it Assertion (1)}~:
Notons que l'idéal de l'anneau local $A_\infty$ est de type
fini. Dès lors, le complété $\widehat{A_\infty}$ est noethérien
d'après \isoegazero I {7.2.7, 7.2.8}.
Il en est de même pour $\widehat{R_\infty}$ et $\widehat{B_\infty}$.

\medskip

\noindent {\it Assertion (2)}~:
D'après le critère de platitude par fibres (cf \isoega 4 {11.3.10}),
 il suffit de montrer que pour tout entier $l>0$,
  $A/\mathfrak{m}_A^l \rightarrow \widehat{A_\infty}/\mathfrak{m}_{\widehat{A_\infty}}^l$.
Or le membre de gauche n'est autre que $(A/\mathfrak{m}_A^l)^{I/\F}$
 qui est plat sur $A/\mathfrak{m}_A^l$ car ce dernier est cohérent
d'après la propriété (v) du lemme \ref{21-lm-ultra-prod}.
Notons que l'extension résiduelle $\kappa_A^{I/\F}/\kappa_A$ 
de $A_\infty/A$ est séparable. En effet,
 pour toute extension finie $L/\kappa_A$,
 $L \otimes_{\kappa_A} \kappa_A^{I/\F}=L^{I/\F}$ est un corps
  d'après le point (i) du lemme \ref{21-lm-ultra-prod}.
On peut alors conclure par application du 
{\it théorème de localisation de la lissité formelle}
d'André (cf \cite{Localisation@Andre}).

\medskip

\noindent {\it Assertion (3)}~:
Pour cette assertion,
il suffit d'appliquer le lemme suivant
à l'idéal maximal de $A_\infty$.
\begin{lemme}
Soit $A$ un anneau et $I$ un idéal de type fini tel que
 $I \subset \mathrm{rad}(A)$. 
Soit $\hat A$ le complété $I$-adique de $A$. \\
Alors le morphisme induit
 $\Spec(\hat A)-V(I.\hat A) \rightarrow \Spec(A)-V(I)$
est surjectif sur les points fermés.
\end{lemme}
Pour ce lemme, on doit montrer que pour tout fermé 
$Z \subset \Spec(A)$,
$$
Z \subset V(I) \Leftrightarrow f^{-1}(Z) \subset V(I.\hat A).
$$
Comme $I$ est de type fini, ceci équivaut à montrer que pour 
tout idéal $J \subset A$,
$$
\exists n>0 \mid I^n \subset J
  \Leftrightarrow
\exists n>0 \mid I^n.\hat A \subset J.\hat A.
$$
Or cela résulte facilement du lemme de Nakayama
puisque pour tout $m>0$, $I^m \subset \mathrm{rad}(A)$.

\medskip

\noindent {\it Assertion (4)}~: 
Remarquons d'abord que 
$\widehat{A_\infty}$ est normal
car $A$ est normal et les fibres géométriques de
 $A \rightarrow  \widehat{A_\infty}$ sont régulières
(cf \isoega 4 {6.14.1}).
Or, $\widehat{B_\infty}/\widehat{R_\infty}$ est libre (de rang $nm$),
donc sans torsion. Comme $\widehat{A_\infty}/\widehat{R_\infty}$
est entière, on en déduit que $\widehat{B_\infty}/\widehat{A_\infty}$
est sans torsion. 
Ainsi, $\widehat{B_\infty}/\widehat{A_\infty}$ est plat.

Or, des relations \eqref{21-disc} pour tout $i \in \NN$,
on déduit
$$
\frac{\disc_{\widehat{B_\infty}/\widehat{R_\infty}}}
 {\disc_{\widehat{A_\infty}/\widehat{R_\infty}}^n} \in \widehat R_\infty^\times
$$
Si $\p$ est un idéal de hauteur $1$ de $\widehat{R_\infty}$,
l'extension $\widehat{B_\infty})_\p/\widehat{A_\infty})_\p$ est sans torsion,
donc libre. 
La relation précédente montre que
$\disc_{\widehat{B_\infty})_\p/\widehat{A_\infty})_\p}$ est inversible
ce qui prouve que $\widehat{B_\infty})_\p/\widehat{A_\infty})_\p$ est 
étale, ce qui démontre (4).

\medskip

\noindent {\it Assertion (5)}~:
D'après le critère de platitude par fibre (cf \isoega 4 {11.3.10}),
pour tout point $x \in \hat X'_\infty$, $P_\infty)_x/X_\infty)_x$ est plat.
Comme dans le point (4), la relation \eqref{21-disc} permet de montrer que
$P_\infty)_x/X_\infty)_x$ est étale. La propriété (5) résulte donc du point 
(3).


\ifx\danslelivre\undefined
\bibliography{../fichiers-configuration-latex/bibliographie}
\bibliographystyle{../fichiers-configuration-latex/smfalpha}
\end{document}
\fi
